\documentclass[10pt,a4paper]{amsart}
\usepackage[margin=19mm]{geometry}
\usepackage{amsmath,amssymb,mathtools}
\usepackage[scr=rsfs,cal=euler]{mathalfa}
\usepackage{xfrac}
\usepackage[unicode,hidelinks,bookmarks=false]{hyperref}
\newcommand{\ie}{i.e.\ }
\newcommand{\cf}{cf.\ }
\newcommand{\resp}{respectively\ }
\newcommand{\Subsection}{\subsection}
\providecommand{\nobreakdash}{\nobreak\hbox{-}}
\setlength{\emergencystretch}{2em}
\multlinegap0pt
\usepackage{amssymb}
\usepackage{mathtools}
\usepackage[shortlabels]{enumitem}
\usepackage{tikz}
\newtheorem{definition}{Definition}
\newtheorem{theorem}{Theorem}
\newcommand{\R}{\mathbb{R}}
\DeclareMathOperator{\dt}{det}
\newcommand{\ip}[2]{\langle #1,\, #2\rangle}
\newcommand{\tr}{\operatorname{tr}}
\title{An intrinsic characterization of the Bogdanov--Takens normal-form coefficients and a mixed-volume obstruction to non-isolated degeneracies}
\author{V\'ictor Castellanos, Ram\'on Eduardo Chan-L\'opez}
\date{Source: arXiv:2608.13931v1 (CC BY 4.0)}
\begin{document}
\maketitle

\noindent\emph{Source and licence.} An intrinsic characterization of the Bogdanov--Takens normal-form coefficients and a mixed-volume obstruction to non-isolated degeneracies. Version arXiv:2608.13931v1 (CC BY 4.0), distributed by arXiv under the Creative Commons Attribution 4.0 International licence, http://creativecommons.org/licenses/by/4.0/, as recorded in the arXiv licence metadata for this exact version and shown on the abstract page https://arxiv.org/abs/2608.13931v1. Full licence notice: \texttt{licenses/arxiv-2608.13931-CC-BY-4.0.txt}.\par\smallskip

\noindent\textit{Editorial scope.} Counted unit: the article's theorem thm:A (source0.tex lines 374--378). The article's complete original proof of it is delivered with it (source0.tex lines 380--411, one \texttt{proof} environment of 32 lines). No mathematics is written, repaired or re-derived: the statement and the proof are the article's own lines, in the article's own notation, and no retained line is substituted or re-flowed.  Retained uncounted as the source-local context the proof needs: lines 136--140; lines 229--245.  Nothing was omitted from any retained span, so the package carries no omission record.  No retained line contains a citation, so the wrapper carries no bibliography and no bibliography entry.  No retained line contains a figure environment, an image-inclusion command or drawing code, so no image is delivered and the package declares no asset dependency.  Editorial formatting only, in addition: an AMS \texttt{amsart} preamble that restates the article's own declarations and macros used by the retained lines, namely the environments \texttt{definition}, \texttt{theorem} on separate counters, together with the standard packages those lines need.  The retained environments print in this document as Definition 1, Theorem 1 in order of appearance, whereas the article prints the same items with the numbering of its own full document.  The complete native source of the article as distributed in the arXiv e-print for this exact version is bundled unchanged as \texttt{sources/207626/source0.tex}.
	\begin{equation}\label{eq:ThmA}
		a=-\tfrac12\,\ip{\nabla\!\dt DX(p)}{q_{0}},
		\qquad
		b=\ip{\nabla\!\tr DX(p)}{q_{0}} .
	\end{equation}

	\begin{definition}\label{def:kbasis}
		A \emph{Kuznetsov basis adapted to $q_{0}$} is a quadruple
		$(q_{0},q_{1},p_{0},p_{1})$ of vectors of $\R^{2}$ with
		\begin{equation}\label{eq:kbasis}
			Jq_{0}=0,\quad Jq_{1}=q_{0},\quad
			J^{\top}p_{1}=0,\quad J^{\top}p_{0}=p_{1},
		\end{equation}
		\begin{equation}\label{eq:biorth}
			\ip{q_{0}}{p_{0}}=\ip{q_{1}}{p_{1}}=1,\qquad
			\ip{q_{0}}{p_{1}}=\ip{q_{1}}{p_{0}}=0 .
		\end{equation}
		The associated BT coefficients are
		\begin{equation}\label{eq:abdef}
			a=\tfrac12\ip{p_{1}}{B(q_{0},q_{0})},\qquad
			b=\ip{p_{0}}{B(q_{0},q_{0})}+\ip{p_{1}}{B(q_{0},q_{1})} .
		\end{equation}
	\end{definition}

	\begin{theorem}[Theorem A]\label{thm:A}
		Let $X$ be of class $C^{2}$ with $X(p)=0$ and $J=DX(p)$ nilpotent of
		rank one, let $q_{0}$ span $\ker J$ and let $a,b$ be as in
		Definition~\ref{def:kbasis}.  Then \eqref{eq:ThmA} holds.
	\end{theorem}

	\begin{proof}
		Work in the adapted basis $(q_{0},q_{1})$ of
		Definition~\ref{def:kbasis}, with dual basis $(p_{0},p_{1})$, so that
		\[
		[J]_{(q_{0},q_{1})}=\begin{pmatrix}0&1\\ 0&0\end{pmatrix}.
		\]
		Let $L:=D(DX)_{p}[q_{0}]$ be the derivative of the matrix-valued map
		$z\mapsto DX(z)$ at $p$ in the direction $q_{0}$; by definition of the
		second-order form, $L(w)=B(q_{0},w)$.  In the same basis,
		\[
		[L]=\begin{pmatrix}
			\ip{p_{0}}{B(q_{0},q_{0})} & \ip{p_{0}}{B(q_{0},q_{1})}\\[2pt]
			\ip{p_{1}}{B(q_{0},q_{0})} & \ip{p_{1}}{B(q_{0},q_{1})}
		\end{pmatrix}.
		\]
		
		Since the trace is linear,
		\[
		d(\tr DX)_{p}(q_{0})=\tr L
		=\ip{p_{0}}{B(q_{0},q_{0})}+\ip{p_{1}}{B(q_{0},q_{1})}=b,
		\]
		which is \eqref{eq:abdef}.  For the determinant, differentiating at the
		matrix $J$ gives $d(\dt)_{J}(L)=\tr\bigl(\operatorname{adj}(J)\,L\bigr)$,
		and
		$\operatorname{adj}\left(\begin{smallmatrix}0&1\\0&0\end{smallmatrix}\right)
		=\left(\begin{smallmatrix}0&-1\\0&0\end{smallmatrix}\right)$, so only
		the $(2,1)$ entry of $[L]$ survives:
		\[
		d(\dt DX)_{p}(q_{0})=-\ip{p_{1}}{B(q_{0},q_{0})}=-2a .
		\]
		Both identities are \eqref{eq:ThmA}.
	\end{proof}

\end{document}
